% Source: https://artofproblemsolving.com/community/c3h1342157_pythagorean_equation
% Stable identity: aops-topic-1342157-q01
\textbf{Proof.} (The three formulas are standard and can be proved by induction on $n$.) Substituting them:
\[
a+6b+3c=T^2+2T^2(4T-1)+T^2\left(6T^2-4T+1\right)=T^2\left(6T^2+4T\right)=2T^3(3T+2),
\]
\[
7a+13b+4c=\dfrac{T^2\left(21+52T-13+24T^2-16T+4\right)}{3}=\dfrac{T^2\left(24T^2+36T+12\right)}{3}=4T^2(2T+1)(T+1),
\]
\[
7a+14b+5c=\dfrac{T^2\left(21+56T-14+30T^2-20T+5\right)}{3}=\dfrac{T^2\left(30T^2+36T+12\right)}{3}=2T^2\left(5T^2+6T+2\right).
\]
Dividing the claimed identity by $\left(2T^2\right)^2$, it becomes
\[
\left(3T^2+2T\right)^2+\left(2(2T+1)(T+1)\right)^2=\left(5T^2+6T+2\right)^2 .
\]
The left side is $9T^4+12T^3+4T^2+4\left(4T^4+12T^3+13T^2+6T+1\right)=25T^4+60T^3+56T^2+24T+4$, and the right side is $25T^4+60T^3+(36+20)T^2+24T+4$, the same polynomial. This proves the identity for every $n$.
